\paragraph{Condition d'émergence d'un rayon lumineux} 
La troisième loi de Snell-Descartes énonce :
\[\sin(r) = \frac{n_1}{n_2} \cdot \sin(i) \]

\noindent Où $n_1$, $n_2 \in \mathbb{R}$ 
désignent les indices optiques des milieux d'étude 
et $i$ et $r \in \mathbb{R}$ désignent respectivement les angles incident et réfléchi.  
\medskip 

\noindent Or, pour tout \(r \in \mathbb{R}\), 
la fonction sinus est une fonction de $\mathbb{R}$ dans $[0, 1]$

\noindent Cela a pour conséquence que $\sin(r) \leq 1$ 
et donc que :
\[\frac{n_1}{n_2} \cdot \sin(i) \leq 1\]

\noindent Alors, il ne peut y avoir existence d'un rayon réfracté uniquement lorsque : 
\[\sin(i) \leq \frac{n_2}{n_1}\]

\noindent On pose $i = \alpha_{lim}$, 
l'expression de l'angle limite de réfraction $\alpha_{lim}$ s'exprime : 
\begin{equation}
    \alpha_{lim} = \sin^{-1}\left (\frac{n_2}{n_1} \right )
    \label{eq:alpha_lim}
\end{equation}

\paragraph{Condition d'émergence du rayon incident en $I$ et $I'$}
En $I'$, 
la condition d'émergence est donnée par :  
\[r' < \arcsin\left(\frac{1}{n}\right)\]
Or, d'après les relations du prisme, 
\[r = A - r'\]
L'émergence impose alors 
\[r > A - \arcsin\left(\frac{1}{n}\right)\]
En appliquant Snell-Descartes en $I$, 
\[\sin(i) = n \sin(r)\]
On en déduit l'incidence minimale 
\[sin(i) > n \sin\left( A - \arcsin\left(\frac{1}{n}\right) \right)\]